% ============================================================================
% 08_mlflow_prompt_engineering.tex — Capítulo 8: Del MLOps al LLMOps:
% Experimentación de Prompt Engineering y Prompt Registry en MLflow 3
% ============================================================================

\chapter[LLMOps, Prompt Engineering y Prompt Registry]{Del MLOps al LLMOps: Experimentación de Prompt Engineering y Prompt Registry en MLflow}\label{chap:mlflow-prompting}

En el Capítulo~\ref{chap:mlflow-clasico} hemos analizado cómo el ecosistema MLOps tradicional ---y en particular \pkg{MLflow}--- disciplina el entrenamiento, el versionado y el despliegue de modelos supervisados clásicos, desde árboles de decisión potenciados por gradiente hasta redes neuronales profundas. En aquel paradigma, el valor intelectual y operativo reside fundamentalmente en la matriz de pesos $\theta \in \mathbb{R}^d$ aprendida a partir de un conjunto de entrenamiento tabular o estructurado, junto con los transformadores de características que garantizan la consistencia entre entrenamiento e inferencia.

La irrupción de los Grandes Modelos de Lenguaje (\textit{Large Language Models}, LLMs) y de los sistemas de Inteligencia Artificial Generativa (GenAI) en entornos productivos entre 2024 y 2026 ha desplazado drásticamente ese centro de gravedad. En la gran mayoría de las aplicaciones empresariales modernas ---asistentes de soporte, extracción estructurada de documentos, generación aumentada por recuperación (RAG) o agentes autónomos--- las organizaciones no entrenan un modelo fundacional desde cero. En su lugar, orquestan modelos preentrenados (consumidos mediante APIs gestionadas o servidos en infraestructura propia con \pkg{vLLM}) cuyo comportamiento efectivo viene determinado por un artefacto de naturaleza completamente distinta: \textbf{las instrucciones en lenguaje natural (\textit{prompts}), los ejemplos en contexto, el conocimiento recuperado dinámicamente y las definiciones de herramientas}.

Este cambio de paradigma inaugura la disciplina del \textbf{LLMOps} (\textit{Large Language Model Operations}). Lejos de quedar obsoleto, \pkg{MLflow 3} ha sido rediseñado desde sus cimientos para actuar como plano unificado tanto para MLOps clásico como para ingeniería GenAI, elevando el \textit{prompt engineering} desde el ensayo artesanal en cuadernos efímeros hasta una disciplina experimental rigurosa, versionada, optimizable algorítmicamente y gobernable mediante CI/CD. En este capítulo diseccionamos esa transformación: primero caracterizamos qué cambia exactamente respecto al MLOps clásico; a continuación diseñamos experimentos reproducibles de \textit{prompt engineering} en MLflow Tracking; después profundizamos en el \textbf{MLflow Prompt Registry} (\code{mlflow.genai}); exploramos la optimización automática mediante \textbf{DSPy} y \textbf{GEPA} analizando la frontera de Pareto entre calidad y coste; y cerramos con la construcción de \textit{Golden Datasets} y puertas de regresión automatizadas.

% ============================================================================
\section[ML Clásico vs. LLMOps y GenAI]{Qué cambia al pasar de ML clásico a LLMs y GenAI}\label{sec:mlops-vs-llmops}
% ============================================================================

Un error recurrente al incorporar LLMs en equipos con experiencia previa en ciencia de datos consiste en asumir que las prácticas de MLOps clásico se transfieren sin fricción, o bien caer en el extremo opuesto y abandonar cualquier rigor experimental al considerar que escribir \textit{prompts} es una actividad meramente lingüística. Ambos enfoques fracasan en producción porque ignoran las diferencias estructurales entre optimizar una función paramétrica diferenciable y programar un sistema estocástico de lenguaje natural.

En el aprendizaje automático supervisado clásico, el ciclo iterativo se organiza alrededor de un algoritmo de optimización numérica que minimiza una función de pérdida bien definida $\mathcal{L}(\theta; \mathcal{D}_{\text{train}})$. Una vez entrenado y serializado el artefacto $\hat{\theta}$, la inferencia sobre un vector de entrada $x$ es esencialmente determinista, su coste computacional por predicción es predecible y acotado, y la evaluación contra un conjunto de validación etiquetado se reduce a calcular métricas cerradas como precisión, \textit{recall}, $F_1$, AUC-ROC o RMSE.

En un sistema basado en LLMs, por el contrario, el modelo fundacional suele permanecer congelado (o actualizado unilateralmente por un proveedor externo) y la ``programación'' del comportamiento se desplaza a un espacio discreto, combinatorio y de altísima sensibilidad semántica: el espacio de los \textit{prompts} y del contexto inyectado en la ventana de entrada. Pequeñas alteraciones aparentemente inocuas ---reordenar dos reglas en el \textit{system prompt}, cambiar un delimitador de Markdown por etiquetas XML o sustituir un ejemplo \textit{few-shot}--- pueden inducir variaciones abruptas de más de 15 puntos porcentuales en la tasa de cumplimiento de un esquema JSON o en la propensión a alucinar. La Tabla~\ref{tab:mlops-vs-llmops} sintetiza las diferencias arquitectónicas y operativas clave entre ambos paradigmas.

\begin{table}[htbp]
\centering
\caption{Comparativa estructural entre MLOps clásico y LLMOps/GenAI en entornos productivos.}
\label{tab:mlops-vs-llmops}
\small
\begin{tabularx}{\textwidth}{@{}L{2.7cm} Y Y@{}}
\toprule
\rowcolor{TableHeaderBg}
\textcolor{white}{\sffamily\bfseries Dimensión} &
\textcolor{white}{\sffamily\bfseries MLOps Clásico (Tabular / DL)} &
\textcolor{white}{\sffamily\bfseries LLMOps / GenAI (LLMs, RAG, Agentes)} \\
\midrule
Artefacto principal versionado &
Matriz de pesos $\hat{\theta}$, hiperparámetros de entrenamiento y \textit{pipeline} de preprocesado de \textit{features}. &
Plantillas de \textit{prompt} (sistema y usuario), ejemplos \textit{few-shot}, esquema de herramientas, configuración RAG e identificador del modelo base. \\
\rowcolor{TableStripe}
Naturaleza de la inferencia &
Determinista para una entrada $x$ fija; salida numérica o categórica acotada. &
Estocástica (controlada por \code{temperature}, \code{top\_p}); incluso con $T=0$ existe variabilidad por paralelismo GPU/MoE. \\
Evaluación y verdad terreno &
Supervisada y cerrada: etiquetas exactas $y_i$ permiten cálculo directo de error sin ambigüedad. &
Respuestas abiertas con múltiples formulaciones válidas; requiere combinación de validadores deterministas, \textit{LLM-as-a-Judge} y revisión humana. \\
\rowcolor{TableStripe}
Métricas críticas de negocio &
$F_1$, AUC-ROC, RMSE, Log-Loss, calibración probabilística y latencia por lote o petición. &
Calidad semántica, tasa de alucinación (\textit{groundedness}), cumplimiento de formato, TTFT, \textit{tokens}/s y coste USD por cada 1\,000 peticiones. \\
Perfil de costes y iteración &
Coste concentrado en entrenamiento intensivo (GPUs/Spark); inferencia unitaria muy barata. &
Iteración rápida sin entrenamiento de pesos; coste dominante en inferencia productiva y en baterías de evaluación con jueces LLM. \\
\rowcolor{TableStripe}
Modos de fallo y deriva &
\textit{Data drift}, \textit{concept drift}, valores nulos inesperados y \textit{training-serving skew}. &
\textit{Prompt drift}, \textit{model drift} silencioso de APIs externas, \textit{context rot}, alucinaciones sutiles, \textit{prompt injection} y fuga de PII. \\
Despliegue y \textit{rollback} &
Reentrenamiento, empaquetado en imagen Docker/artefacto y redespliegue del servicio de inferencia. &
Desacoplado: promoción o reversión instantánea del alias \code{@production} en Prompt Registry sin recompilar el microservicio. \\
\bottomrule
\end{tabularx}
\end{table}

\subsection{Por qué el \textit{prompt engineering} artesanal en notebooks fracasa}\label{subsec:fracaso-artesanal}

Cuando un prototipo generativo pasa de una demostración interna a un servicio con miles de usuarios diarios, el flujo de trabajo artesanal ---caracterizado por cadenas multilínea incrustadas directamente en el código fuente o probadas manualmente en un cuaderno interactivo--- colapsa rápidamente por cinco motivos estructurales:

\begin{enumerate}
  \item \textbf{Evaluación por sensaciones (\textit{vibe-checking}).} El desarrollador modifica una frase del \textit{prompt}, comprueba manualmente tres o cuatro ejemplos que habían fallado antes y asume que la nueva versión es mejor. Sin una batería de evaluación fija, es habitual introducir regresiones invisibles en otros segmentos del dominio (efecto \textit{whack-a-mole}: arreglar un caso borde rompe silenciosamente diez casos frecuentes).
  \item \textbf{Acoplamiento entre código de aplicación e instrucciones lingüísticas.} Incrustar los \textit{prompts} como literales de texto dentro de archivos Python obliga a abrir un \textit{Pull Request}, ejecutar pruebas unitarias, construir imágenes de contenedor y redesplegar pods de Kubernetes cada vez que un experto de dominio necesita ajustar una política de tono o una regla de negocio.
  \item \textbf{Pérdida total de reproducibilidad y trazabilidad.} Meses después de un despliegue, nadie recuerda con qué combinación exacta de versión de \textit{prompt}, temperatura, versión del modelo base y estrategia de recuperación se obtuvieron las métricas presentadas al comité de producto.
  \item \textbf{Ceguera económica y operacional.} Un \textit{prompt} enriquecido con seis ejemplos largos y cadena de pensamiento puede elevar la exactitud del 88\,\% al 90\,\%, pero triplicar el número de \textit{tokens} de entrada y duplicar la latencia p95. Sin instrumentación experimental homogénea, los equipos optimizan una única dimensión a ciegas y descubren la factura al final del mes.
  \item \textbf{Fragilidad ante cambios del modelo subyacente.} Los proveedores de APIs retiran o actualizan periódicamente sus modelos. Un \textit{prompt} sobreajustado manualmente a las peculiaridades de un modelo concreto (por ejemplo, \code{gpt-4o}) puede degradarse severamente al migrar a otro proveedor (como Claude o Gemini) o al destilar la carga a un modelo más pequeño y económico.
\end{enumerate}

\begin{warningbox}[El mito de \code{temperature=0} y el \textit{model drift} silencioso]
En LLMOps existen dos fuentes de no reproducibilidad que suelen sorprender a los ingenieros procedentes del ML clásico:
\begin{itemize}
  \item \textbf{No determinismo residual con \code{temperature=0}:} incluso fijando la temperatura a cero y proporcionando un parámetro \code{seed}, los motores de inferencia modernos con arquitecturas \textit{Mixture-of-Experts} (MoE) y \textit{continuous batching} presentan pequeñas variaciones en el orden de suma en coma flotante en GPU. Ante \textit{logits} casi empatados entre dos \textit{tokens}, esto provoca respuestas distintas para exactamente la misma petición. Por ello, los experimentos rigurosos deben reportar medias y desviaciones sobre conjuntos amplios y no confiar en una única ejecución aislada.
  \item \textbf{\textit{Model drift} silencioso en alias flotantes:} apuntar en producción a nombres genéricos como \code{gpt-4o} o \code{claude-sonnet} delega en el proveedor el momento en que cambian los pesos subyacentes. En sistemas críticos, tanto el código productivo como los registros de MLflow deben fijar identificadores inmutables de instantánea (\textit{model snapshot}, por ejemplo \code{gpt-4o-2024-11-20}) y registrar cualquier migración como un experimento controlado.
\end{itemize}
\end{warningbox}

Frente a estas limitaciones, \pkg{MLflow 3} reconvierte el \textit{prompt engineering} en ingeniería experimental clásica pero adaptada a la semántica de los LLMs: cada combinación de plantilla, hiperparámetros de decodificación, modelo base y configuración de contexto se ejecuta de forma sistemática sobre un conjunto de evaluación versionado, registrando de manera atómica sus parámetros, métricas de calidad, costes, latencias y artefactos completos de entrada/salida. A su vez, el \textbf{Prompt Registry} desacopla la evolución de las instrucciones del ciclo de despliegue del código, aportando versionado inmutable, linaje y promoción gobernada por alias.

% ============================================================================
\section[Experimentos de Prompting en Tracking]{Diseño de experimentos de Prompt Engineering en MLflow Tracking}\label{sec:exp-prompt-tracking}
% ============================================================================

Tratar los \textit{prompts} como hipótesis científicas verificables exige estructurar dos elementos antes de escribir una sola línea de código: (i) una taxonomía clara de las técnicas de diseño de \textit{prompts} que conforman el espacio de búsqueda, y (ii) un esquema normalizado de parámetros, métricas y artefactos en \pkg{MLflow Tracking} que permita comparar ejecuciones heterogéneas en igualdad de condiciones.

\subsection{Taxonomía de técnicas evaluables sistemáticamente}\label{subsec:taxonomia-prompting}

En lugar de iterar frases arbitrarias, un equipo maduro de LLMOps parametriza sus experimentos recorriendo familias bien definidas de técnicas de \textit{prompt engineering}, evaluando su impacto marginal en calidad, latencia y consumo de \textit{tokens}:

\begin{description}
  \item[Zero-Shot estructurado] Proporciona únicamente la definición de la tarea, restricciones explícitas, delimitadores claros (por ejemplo etiquetas XML \code{<documento>...</documento>} para separar instrucciones de datos no confiables) y el formato de salida esperado, sin incluir ejemplos resueltos. Constituye siempre la línea base (\textit{baseline}) obligatoria por su bajo coste en \textit{tokens} de entrada.
  \item[Few-Shot estático vs. dinámico] Incorpora $k$ pares entrada--salida ejemplares dentro del contexto. En la variante \textbf{estática}, los $k$ ejemplos son fijos y se seleccionan manualmente para cubrir casos representativos y casos borde. En la variante \textbf{dinámica}, cada consulta entrante se proyecta a un espacio vectorial de \textit{embeddings} y se recuperan dinámicamente los $k$ ejemplos más similares semánticamente de un banco curado. El \textit{few-shot} dinámico suele superar claramente al estático en dominios con muchas categorías, pero introduce latencia adicional de recuperación e invalida parcialmente la caché de prefijo (\textit{prompt caching}) del proveedor LLM.
  \item[Chain-of-Thought (CoT) y razonamiento guiado] Instruye al modelo a explicitar pasos intermedios de verificación antes de emitir la conclusión final. Resulta decisivo en tareas con reglas de negocio condicionales, cálculos aritméticos o clasificación con múltiples criterios excluyentes. Su contrapartida directa es el incremento de \textit{tokens} de salida (\textit{completion tokens}), que son entre $3\times$ y $5\times$ más caros que los de entrada y penalizan la latencia total. Con modelos razonadores nativos de última generación (que ya generan razonamiento interno oculto), un CoT verboso en el \textit{prompt} puede ser redundante o incluso contraproducente, lo que refuerza la necesidad de medirlo empíricamente.
  \item[Structured Outputs (JSON Schema / Pydantic)] Restringe la decodificación del modelo mediante gramáticas libres de contexto (\textit{constrained decoding}) derivadas de un esquema JSON o de un modelo \pkg{Pydantic}. Elimina prácticamente los errores de sintaxis en flujos de extracción de entidades, enrutamiento y llamadas a APIs, aunque conviene recordar que garantiza validez \emph{sintáctica}, no veracidad \emph{semántica} (un campo puede respetar el tipo \code{float} y contener un importe inventado).
  \item[Role \& Persona Prompting con políticas negativas] Asigna un rol especializado y define explícitamente qué debe hacer el modelo ante incertidumbre o ausencia de evidencia (por ejemplo, devolver \code{"INSUFICIENTE\_EVIDENCIA"} en lugar de conjeturar). Las cláusulas explícitas de abstención son una de las palancas más eficaces para reducir alucinaciones en aplicaciones empresariales.
  \item[Descomposición de tareas y patrón ReAct] Divide un problema complejo en subtareas secuenciales (encadenamiento de \textit{prompts} especializados) o alterna ciclos explícitos de razonamiento (\textit{Thought}), invocación de herramientas (\textit{Action}) y observación de resultados (\textit{Observation}). Reduce la carga cognitiva por llamada, pero multiplica el número de rondas de red y exige trazabilidad distribuida (tema central del Capítulo~\ref{chap:mlflow-llms}).
\end{description}

\begin{table}[htbp]
\centering
\caption{Taxonomía de técnicas de \textit{Prompt Engineering} y sus compromisos operativos.}
\label{tab:taxonomia-prompting}
\small
\begin{tabularx}{\textwidth}{@{}L{2.9cm} Y C{2.1cm} C{2.1cm} L{2.8cm}@{}}
\toprule
\rowcolor{TableHeaderBg}
\textcolor{white}{\sffamily\bfseries Técnica} &
\textcolor{white}{\sffamily\bfseries Caso de uso ideal} &
\textcolor{white}{\sffamily\bfseries Impacto en \textit{tokens}} &
\textcolor{white}{\sffamily\bfseries Impacto en latencia} &
\textcolor{white}{\sffamily\bfseries Riesgo principal} \\
\midrule
Zero-Shot estructurado &
Tareas bien definidas, resumen, clasificación simple y \textit{baseline}. &
Bajo &
Bajo &
Ambigüedad en casos borde y formato inconsistente. \\
\rowcolor{TableStripe}
Few-Shot estático &
Estilos específicos, taxonomías cerradas y extracción homogénea. &
Medio (cacheable) &
Bajo--Medio &
Sesgo hacia las clases sobrerrepresentadas en los ejemplos. \\
Few-Shot dinámico &
Dominios amplios con decenas de categorías o reglas dispares. &
Medio (no cacheable) &
Medio &
Ejemplos recuperados ruidosos que confunden al modelo. \\
\rowcolor{TableStripe}
Chain-of-Thought &
Lógica condicional compleja, auditoría de políticas y cálculo. &
Alto (salida) &
Alto &
Racionalización \textit{post-hoc} y sobrecoste en modelos razonadores. \\
Structured Outputs &
Extracción hacia bases de datos, enrutadores y \textit{pipelines} ETL. &
Bajo--Medio &
Bajo &
Validez sintáctica sin veracidad semántica (alucinación tipada). \\
\rowcolor{TableStripe}
Descomposición / ReAct &
Consultas multi-paso con acceso a APIs, SQL o búsqueda documental. &
Muy alto &
Muy alto &
Bucles infinitos de herramientas y propagación de errores. \\
\bottomrule
\end{tabularx}
\end{table}

\subsection{Variables, métricas y economía experimental en MLflow Tracking}\label{subsec:variables-metricas-tracking}

Al igual que en el ajuste de hiperparámetros de un algoritmo \pkg{XGBoost} se registran \code{max\_depth} y \code{learning\_rate}, un experimento de \textit{prompt engineering} en \pkg{MLflow Tracking} debe capturar sistemáticamente tres planos complementarios: \textbf{configuración experimental} (\code{log\_params}), \textbf{métricas multidimensionales} (\code{log\_metrics}) y \textbf{evidencia granular por ejemplo} (\code{log\_table} y \code{log\_dict}).

\textbf{Hiperparámetros y metadatos de inferencia (\code{log\_params}).} Todo \textit{run} debe registrar de forma inmutable las variables que definen la configuración evaluada:
\begin{itemize}
  \item \textbf{Modelo e inferencia:} \code{model\_name}, \code{model\_snapshot}, \code{temperature}, \code{top\_p}, \code{max\_tokens}, \code{seed} y \code{response\_format\_schema}.
  \item \textbf{Identidad del \textit{prompt}:} \code{prompt\_name}, \code{prompt\_version}, \code{prompt\_strategy} (por ejemplo \code{cot\_structured}), \code{system\_prompt\_hash} (resumen criptográfico SHA-256 del texto exacto para detectar cambios de un solo carácter) y \code{few\_shot\_k}.
  \item \textbf{Contexto RAG (si aplica):} \code{embedding\_model}, \code{chunk\_size}, \code{chunk\_overlap}, \code{retriever\_top\_k} y \code{reranker\_enabled}.
  \item \textbf{Conjunto de evaluación:} \code{eval\_dataset\_name}, \code{eval\_dataset\_version} y \code{eval\_sample\_count}.
\end{itemize}

\textbf{Métricas de calidad, operación y economía (\code{log\_metrics}).} Evaluar una plantilla solo por exactitud produce decisiones arquitectónicas erróneas. Cada \textit{run} debe consolidar métricas en tres ejes:
\begin{enumerate}
  \item \textbf{Calidad funcional y semántica:} tasa de validez sintáctica (\code{json\_valid\_rate}), coincidencia exacta en campos estructurados (\code{exact\_match}), $F_1$ a nivel de \textit{tokens} o entidades (\code{token\_f1}), puntuaciones de jueces LLM (\code{judge\_correctness}, \code{groundedness\_score}) y tasa de abstención correcta en casos fuera de dominio (\code{abstention\_accuracy}).
  \item \textbf{Rendimiento operacional:} mediana y cola de latencia total (\code{latency\_ms\_p50}, \code{latency\_ms\_p95}), tiempo hasta el primer \textit{token} (\code{ttft\_ms\_p50}, crítico en interfaces conversacionales en \textit{streaming}) y velocidad de generación (\code{tokens\_per\_second}).
  \item \textbf{Economía unitaria:} consumo acumulado y medio de \textit{tokens} de entrada y salida (\code{prompt\_tokens\_avg}, \code{completion\_tokens\_avg}), coste total del experimento (\code{estimated\_cost\_usd}) y coste normalizado por cada mil peticiones (\code{cost\_per\_1k\_requests\_usd}).
\end{enumerate}

Formalmente, dado un conjunto de evaluación de $N$ consultas y un modelo $m$ con tarifas por millón de \textit{tokens} de entrada $P_{\text{in}}(m)$ y de salida $P_{\text{out}}(m)$ (expresadas en USD), el coste estimado del experimento y el coste unitario escalado a producción vienen dados por:
\begin{align}\label{eq:coste-llm}
  \text{Coste}_{\text{run}}(m, p) &= \sum_{i=1}^{N} \left( \frac{T^{\text{in}}_{i}(p)}{10^{6}} \cdot P_{\text{in}}(m) + \frac{T^{\text{out}}_{i}(m, p)}{10^{6}} \cdot P_{\text{out}}(m) \right), \\
  C_{1\text{k}}(m, p) &= \frac{1000}{N} \, \text{Coste}_{\text{run}}(m, p). \nonumber
\end{align}
Cuando el \textit{prompt} de sistema y los ejemplos \textit{few-shot} estáticos superan el umbral mínimo de prefijo del proveedor (habitualmente 1\,024 \textit{tokens}), los mecanismos de \textit{prompt caching} aplican descuentos del 50\,\% al 90\,\% sobre la porción de prefijo repetida $T^{\text{cached}}_{i}$; modelar explícitamente $T^{\text{cached}}_{i}$ dentro de la ecuación~\eqref{eq:coste-llm} evita sobrestimar el coste real de estrategias \textit{few-shot} estáticas frente a las dinámicas.

\subsection{Caso práctico: barrido sistemático de plantillas, temperatura y modelos}\label{subsec:ejemplo-sweep}

\begin{examplebox}[Barrido factorial de triaje de incidencias en MLflow Tracking]
Supongamos un servicio corporativo de atención al cliente que clasifica incidencias entrantes y extrae un objeto JSON estructurado con tres campos: \code{categoria} (\code{facturacion}, \code{tecnico}, \code{envio}, \code{cancelacion}), \code{prioridad} (\code{alta}, \code{media}, \code{baja}) y \code{requiere\_escalado} (booleano). Queremos comparar de forma reproducible \textbf{3 estrategias de \textit{prompt}} (\code{zero\_shot}, \code{few\_shot\_static}, \code{cot\_structured}) $\times$ \textbf{2 temperaturas} (\code{0.0}, \code{0.3}) $\times$ \textbf{2 modelos} (\code{gpt-4o-mini}, \code{gpt-4o}), generando $3 \times 2 \times 2 = 12$ ejecuciones anidadas bajo un \textit{run} padre en MLflow Tracking.
\end{examplebox}

El Listado~\ref{lst:prompt-sweep} presenta la implementación completa del barrido experimental en Python. Obsérvese cómo se utiliza un \textit{run} padre para agrupar toda la campaña experimental, \textit{runs} hijos para cada celda de la matriz, validación estricta mediante \pkg{Pydantic}, cálculo de percentiles de latencia y coste según la ecuación~\eqref{eq:coste-llm}, y registro de resultados fila por fila con \code{mlflow.log\_table} para permitir la inspección comparativa directa en la interfaz gráfica de MLflow.

\begin{lstlisting}[style=python, caption={Barrido sistemático de 3 plantillas $\times$ 2 temperaturas $\times$ 2 modelos registrado en MLflow Tracking.}, label={lst:prompt-sweep}]
import hashlib
import json
import time
from itertools import product
from typing import Literal

import mlflow
import numpy as np
from openai import OpenAI
from pydantic import BaseModel, ValidationError

# 1. Esquema objetivo para validacion determinista de salidas estructuradas
class TicketTriage(BaseModel):
    categoria: Literal["facturacion", "tecnico", "envio", "cancelacion"]
    prioridad: Literal["alta", "media", "baja"]
    requiere_escalado: bool

# 2. Tarifas de referencia por millon de tokens (USD)
PRICING_USD_PER_1M = {
    "gpt-4o-mini-2024-07-18": {"input": 0.15, "output": 0.60},
    "gpt-4o-2024-11-20": {"input": 2.50, "output": 10.00},
}

# 3. Plantillas candidatas bajo estudio
PROMPT_TEMPLATES = {
    "zero_shot": (
        "Eres un clasificador de tickets de soporte. Devuelve SOLO un objeto JSON "
        "valido con las claves: categoria (facturacion|tecnico|envio|cancelacion), "
        "prioridad (alta|media|baja) y requiere_escalado (true|false)."
    ),
    "few_shot_static": (
        "Eres un clasificador de tickets de soporte. Devuelve SOLO JSON valido.\n"
        "Ejemplo 1:\n"
        "Ticket: 'Me han cobrado dos veces la factura mensual.'\n"
        'Salida: {"categoria":"facturacion","prioridad":"alta","requiere_escalado":true}\n'
        "Ejemplo 2:\n"
        "Ticket: 'No puedo descargar el manual PDF desde la app.'\n"
        'Salida: {"categoria":"tecnico","prioridad":"baja","requiere_escalado":false}\n'
        "Ejemplo 3:\n"
        "Ticket: 'Quiero dar de baja mi suscripcion anual hoy mismo.'\n"
        'Salida: {"categoria":"cancelacion","prioridad":"alta","requiere_escalado":true}'
    ),
    "cot_structured": (
        "Eres un especialista senior en triaje de incidencias. Analiza paso a paso "
        "en silencio: (1) identifica el dominio principal del problema, (2) evalua "
        "impacto economico o bloqueo operativo para fijar prioridad, (3) marca "
        "requiere_escalado=true solo si hay cobro indebido, baja inminente o caida "
        "critica. Responde UNICAMENTE con un objeto JSON que cumpla el esquema: "
        '{"categoria": "...", "prioridad": "...", "requiere_escalado": true/false}.'
    ),
}

# 4. Muestra ilustrativa del conjunto de evaluacion etiquetado
EVAL_DATASET = [
    {
        "id": "T-001",
        "texto": "La pasarela de pago rechazo la tarjeta pero el cargo aparece duplicado.",
        "expected": {"categoria": "facturacion", "prioridad": "alta", "requiere_escalado": True},
    },
    {
        "id": "T-002",
        "texto": "El paquete figura entregado pero no ha llegado a mi domicilio.",
        "expected": {"categoria": "envio", "prioridad": "alta", "requiere_escalado": False},
    },
    {
        "id": "T-003",
        "texto": "El boton de exportar CSV no responde en Safari tras la actualizacion.",
        "expected": {"categoria": "tecnico", "prioridad": "media", "requiere_escalado": False},
    },
    {
        "id": "T-004",
        "texto": "Solicito cancelar la cuenta empresarial y reembolso proporcional.",
        "expected": {"categoria": "cancelacion", "prioridad": "alta", "requiere_escalado": True},
    },
]

client = OpenAI()
mlflow.set_experiment("prompt-engineering-triaje-tickets")

def evaluate_combination(strategy: str, system_prompt: str, model: str, temp: float):
    prompt_hash = hashlib.sha256(system_prompt.encode("utf-8")).hexdigest()[:12]
    latencies_ms, ttfts_ms, tps_list = [], [], []
    prompt_tokens_total, completion_tokens_total = 0, 0
    valid_json_count, exact_match_count, field_hits, field_total = 0, 0, 0, 0
    rows_for_table = []

    with mlflow.start_run(run_name=f"{strategy}__{model}__T{temp}", nested=True):
        mlflow.log_params({
            "prompt_strategy": strategy,
            "system_prompt_hash": prompt_hash,
            "model_snapshot": model,
            "temperature": temp,
            "top_p": 1.0,
            "max_tokens": 256,
            "few_shot_k": 3 if strategy == "few_shot_static" else 0,
            "eval_sample_count": len(EVAL_DATASET),
        })

        for sample in EVAL_DATASET:
            t0 = time.perf_counter()
            ttft_ms = None
            chunks = []

            # Streaming habilitado para medir TTFT real y tokens/s
            stream = client.chat.completions.create(
                model=model,
                temperature=temp,
                max_tokens=256,
                response_format={"type": "json_object"},
                stream=True,
                stream_options={"include_usage": True},
                messages=[
                    {"role": "system", "content": system_prompt},
                    {"role": "user", "content": f"<ticket>{sample['texto']}</ticket>"},
                ],
            )
            usage = None
            for event in stream:
                if event.choices and event.choices[0].delta.content:
                    if ttft_ms is None:
                        ttft_ms = (time.perf_counter() - t0) * 1000.0
                    chunks.append(event.choices[0].delta.content)
                if getattr(event, "usage", None):
                    usage = event.usage

            elapsed_ms = (time.perf_counter() - t0) * 1000.0
            raw_output = "".join(chunks).strip()
            latencies_ms.append(elapsed_ms)
            ttfts_ms.append(ttft_ms or elapsed_ms)

            in_tok = usage.prompt_tokens if usage else 0
            out_tok = usage.completion_tokens if usage else 0
            prompt_tokens_total += in_tok
            completion_tokens_total += out_tok
            if elapsed_ms > 0 and out_tok > 0:
                tps_list.append(out_tok / (elapsed_ms / 1000.0))

            # Validacion de esquema y exactitud por campo
            is_valid, is_exact = False, False
            parsed_dict = {}
            try:
                parsed = TicketTriage.model_validate_json(raw_output)
                parsed_dict = parsed.model_dump()
                is_valid = True
                valid_json_count += 1
                expected = sample["expected"]
                hits = sum(parsed_dict[k] == expected[k] for k in expected)
                field_hits += hits
                field_total += len(expected)
                if hits == len(expected):
                    is_exact = True
                    exact_match_count += 1
            except (ValidationError, json.JSONDecodeError):
                field_total += len(sample["expected"])

            rows_for_table.append({
                "id": sample["id"],
                "input_text": sample["texto"],
                "expected_json": json.dumps(sample["expected"], ensure_ascii=False),
                "predicted_raw": raw_output,
                "json_valid": is_valid,
                "exact_match": is_exact,
                "latency_ms": round(elapsed_ms, 1),
                "ttft_ms": round(ttft_ms or elapsed_ms, 1),
                "prompt_tokens": in_tok,
                "completion_tokens": out_tok,
            })

        prices = PRICING_USD_PER_1M[model]
        total_cost = (
            (prompt_tokens_total / 1e6) * prices["input"]
            + (completion_tokens_total / 1e6) * prices["output"]
        )
        n = len(EVAL_DATASET)
        mlflow.log_metrics({
            "json_valid_rate": valid_json_count / n,
            "exact_match": exact_match_count / n,
            "field_accuracy": field_hits / max(field_total, 1),
            "latency_ms_p50": float(np.percentile(latencies_ms, 50)),
            "latency_ms_p95": float(np.percentile(latencies_ms, 95)),
            "ttft_ms_p50": float(np.percentile(ttfts_ms, 50)),
            "tokens_per_second": float(np.mean(tps_list)) if tps_list else 0.0,
            "prompt_tokens_total": prompt_tokens_total,
            "completion_tokens_total": completion_tokens_total,
            "estimated_cost_usd": round(total_cost, 6),
            "cost_per_1k_requests_usd": round((total_cost / n) * 1000.0, 4),
        })

        # Tabla interactiva para comparacion lado a lado en la UI de MLflow
        mlflow.log_table(data=rows_for_table, artifact_file="eval_results_table.json")
        mlflow.log_dict(
            {"system_prompt": system_prompt, "predictions": rows_for_table},
            artifact_file="raw_outputs_audit.json",
        )

# 5. Ejecucion completa del barrido factorial bajo un run padre
with mlflow.start_run(run_name="sweep_triaje_v1"):
    mlflow.set_tags({"campana": "triaje-q4", "tipo": "prompt_grid_search"})
    for (strategy, prompt_text), model, temp in product(
        PROMPT_TEMPLATES.items(),
        list(PRICING_USD_PER_1M.keys()),
        [0.0, 0.3],
    ):
        evaluate_combination(strategy, prompt_text, model, temp)
\end{lstlisting}

Una vez completado el barrido, la interfaz de \pkg{MLflow Tracking} permite ordenar las doce ejecuciones por \code{exact\_match} y \code{cost\_per\_1k\_requests\_usd}, construir gráficos de dispersión directamente en la vista de comparación y abrir la pestaña \textit{Artifacts / Evaluation} para contrastar lado a lado las salidas generadas por \code{eval\_results\_table.json} sobre exactamente los mismos tickets. Este procedimiento sustituye la intuición subjetiva por evidencia auditable.

% ============================================================================
\section[MLflow Prompt Registry]{MLflow Prompt Registry (\texorpdfstring{\code{mlflow.genai}}{mlflow.genai}): Prompts como artefactos de primera clase}\label{sec:prompt-registry}
% ============================================================================

Registrar los resultados de un barrido en MLflow Tracking resuelve el problema de la experimentación, pero deja abierta una cuestión crítica para entornos productivos: \textbf{¿dónde viven, cómo se versionan y cómo se despliegan las plantillas de \textit{prompt} que consumen los microservicios y agentes en producción?}

Hasta 2024, los equipos solían recurrir a soluciones improvisadas: almacenar plantillas en archivos YAML dentro del repositorio Git de la aplicación (acoplando cada cambio lingüístico al ciclo de compilación y despliegue del servicio), guardarlas en tablas ad hoc de PostgreSQL sin trazabilidad con los experimentos, o mantener ficheros dispersos en buckets S3 sin control de versiones ni auditoría de cambios. \pkg{MLflow 3} cierra esta brecha elevando los \textit{prompts} a ciudadanos de primera clase mediante el módulo \code{mlflow.genai} y su \textbf{Prompt Registry}.

\begin{modernbox}[{Arquitectura del Prompt Registry en MLflow 3 \badge[PrimaryAccent]{MLflow 3}}]
El \textbf{MLflow Prompt Registry} replica los principios que convirtieron al \textit{Model Registry} en el estándar de MLOps, pero diseñados específicamente para artefactos de texto parametrizado:
\begin{itemize}
  \item \textbf{Versionado inmutable estilo Git:} cada registro bajo un nombre lógico (por ejemplo \code{soporte\_cliente}) crea una versión numérica inmutable (\code{v1}, \code{v2}, \dots, \code{vN}) acompañada de un mensaje de \textit{commit} (\code{commit\_message}), marca temporal, autor y metadatos clave-valor (\code{tags}). Ninguna versión histórica puede modificarse silenciosamente, garantizando auditoría completa.
  \item \textbf{Plantillas con variables de doble llave:} las variables dinámicas se declaran mediante \code{\{\{variable\}\}}, lo que permite incrustar cómodamente esquemas JSON, fragmentos de código o diccionarios Python dentro del \textit{prompt} sin colisiones de sintaxis.
  \item \textbf{Aliases mutables para gobierno de entornos:} punteros simbólicos como \code{@production}, \code{@staging}, \code{@canary} o \code{@experiment} apuntan a una versión inmutable concreta. Promover un \textit{prompt} o ejecutar un \textit{rollback} instantáneo consiste únicamente en reasignar el alias.
  \item \textbf{Desacoplamiento total código--prompt:} la aplicación consulta la URI \code{prompts:/soporte\_cliente@production} con caché local configurable (TTL en segundos), absorbiendo actualizaciones de negocio sin reconstruir contenedores ni reiniciar pods.
  \item \textbf{Linaje bidireccional con \code{LoggedModel} y evaluaciones:} MLflow vincula automáticamente qué versiones de \textit{prompt} utiliza un modelo/agente registrado y qué ejecuciones de evaluación respaldaron su promoción.
\end{itemize}
\end{modernbox}

La Figura~\ref{fig:prompt-registry-lifecycle} ilustra el ciclo de vida completo de un \textit{prompt} gobernado por el MLflow Prompt Registry, desde la autoría inicial hasta la carga dinámica en producción y el mecanismo de reversión inmediata ante degradaciones detectadas en tiempo real.

\begin{figure}[htbp]
\centering
\resizebox{\textwidth}{!}{%
\begin{tikzpicture}[node distance=1.0cm and 1.45cm]
  % Fila superior: flujo principal de creacion y promocion
  \node[archbox, text width=2.5cm] (auth) {1. Autoría\\[-1pt]{\scriptsize\normalfont Plantilla con \texttt{\{\{var\}\}}}};
  \node[tealbox, text width=2.6cm, right=of auth] (reg) {2. Registro v1\\[-1pt]{\scriptsize\normalfont Inmutable + commit\\+ tags en Registry}};
  \node[archbox, text width=3.0cm, right=of reg] (eval) {3. Evaluación y Trazas\\[-1pt]{\scriptsize\normalfont \texttt{mlflow.genai.evaluate}\\sobre Golden Dataset}};
  \node[amberbox, text width=2.7cm, right=of eval] (opt) {4. Iteración v2..vN\\[-1pt]{\scriptsize\normalfont Ajuste manual o\\optimización GEPA/DSPy}};

  % Fila inferior: puerta de calidad, promocion, produccion y rollback
  \node[darkbox, text width=2.7cm, below=1.55cm of opt] (gate) {5. Quality Gate\\[-1pt]{\scriptsize\normalfont Umbrales calidad,\\seguridad y coste}};
  \node[tealbox, text width=3.0cm, left=of gate] (alias) {6. Promoción Alias\\[-1pt]{\scriptsize\normalfont \texttt{@staging} $\rightarrow$ \texttt{@canary}\\$\rightarrow$ \texttt{@production}}};
  \node[archbox, text width=2.9cm, left=of alias] (prod) {7. Carga Dinámica\\[-1pt]{\tiny\normalfont \texttt{prompts:/...@production}}\\[-1pt]{\scriptsize\normalfont con caché local TTL}};
  \node[amberbox, text width=2.5cm, left=of prod] (mon) {8. Monitorización\\[-1pt]{\scriptsize\normalfont Trazas productivas\\y alertas de deriva}};

  % Recuadro agrupador del Registry (con espacio superior generoso para el bucle interno)
  \coordinate (regtop) at ([yshift=0.95cm]eval.north);
  \begin{scope}[on background layer]
    \node[groupbox, inner sep=12pt, fit=(reg) (opt) (alias) (gate) (regtop)] (reggrp) {};
    \node[anchor=south west, font=\sffamily\scriptsize\bfseries, text=TealAccent]
      at ([xshift=4pt, yshift=2pt]reggrp.north west) {Plano de Gobierno --- MLflow Prompt Registry};
  \end{scope}

  % Flechas principales
  \draw[flowarrow] (auth) -- (reg);
  \draw[flowarrow] (reg) -- (eval);
  \draw[flowarrow] (eval) -- (opt);
  \draw[flowarrow] (opt) -- node[elabel, right=3pt] {Candidata vN} (gate);
  \draw[tealarrow] (gate) -- node[elabel, above=3pt] {Aprobada} (alias);
  \draw[tealarrow] (alias) -- node[elabel, above=3pt] {Resolución URI} (prod);
  \draw[flowarrow] (prod) -- (mon);

  % Bucle de iteracion cuando falla la puerta
  \draw[dashedarrow] (gate.east) -- ++(0.65,0) |- node[elabel, pos=0.25, right=3pt] {Rechazada} (opt.east);

  % Bucle de reevaluacion dentro del Registry (totalmente separado de la etiqueta superior)
  \draw[dashedarrow] (opt.north) -- ++(0,0.62) -| node[elabel, pos=0.25, above=2pt] {Reevaluación continua} (eval.north);

  % Bucle de rollback instantaneo
  \draw[->, >=Stealth, line width=1.1pt, color=CrimsonAccent]
    (mon.south) -- ++(0,-0.70) -| node[elabel, pos=0.25, below=3pt, text=CrimsonAccent] {\bfseries Rollback instantáneo si hay regresión: reasignar \texttt{@production} a v(N--1)} (alias.south);
\end{tikzpicture}%
}
\caption{Ciclo de vida completo de un \textit{prompt} en el MLflow Prompt Registry, desacoplando la iteración y reversión instantánea del ciclo de despliegue de código.}
\label{fig:prompt-registry-lifecycle}
\end{figure}

\subsection{Registro, formateo seguro, gestión de aliases y carga dinámica}\label{subsec:codigo-prompt-registry}

\begin{infobox}[Por qué MLflow utiliza delimitadores de doble llave \code{\{\{variable\}\}}]
En Python estándar, las plantillas basadas en \code{str.format()} emplean llaves simples \code{\{variable\}}. Sin embargo, en ingeniería de \textit{prompts} es sumamente frecuente incluir ejemplos de respuestas JSON o esquemas de herramientas dentro del propio texto del \textit{prompt}, lo que obliga a escapar manualmente decenas de llaves (\code{\{\{} y \code{\}\}}) provocando errores de ejecución sutiles. La sintaxis de doble llave \code{\{\{variable\}\}} adoptada por \code{mlflow.genai} (inspirada en Jinja2/Mustache) elimina esa ambigüedad: las llaves simples de un objeto JSON literal se conservan intactas y solo los identificadores rodeados de dobles llaves se sustituyen de forma segura al invocar \code{prompt.format(...)}.
\end{infobox}

El Listado~\ref{lst:prompt-registry-workflow} muestra el flujo operativo completo con la API de MLflow 3: creación de la versión inicial, iteración a una segunda versión con políticas de citas, promoción mediante aliases (\code{@staging} y \code{@production}), consulta del historial mediante \code{MlflowClient}, carga en producción con caché local y vinculación explícita de linaje con un modelo registrado (\code{LoggedModel}).

\begin{lstlisting}[style=python, caption={Flujo completo de gestión de prompts con MLflow Prompt Registry (\code{mlflow.genai}): registro, formateo, aliases, carga con caché y linaje.}, label={lst:prompt-registry-workflow}]
import mlflow
from mlflow.tracking import MlflowClient
from openai import OpenAI

mlflow.set_tracking_uri("http://mlflow-server.internal:5000")
mlflow.set_experiment("asistente-soporte-produccion")

PROMPT_NAME = "soporte_cliente"

# 1. Registro de la version inicial (v1) del prompt
template_v1 = """Eres el asistente virtual oficial de soporte tecnico.
Responde a la consulta del usuario usando exclusivamente el contexto proporcionado.

<contexto>
{{context}}
</contexto>

<consulta>
{{question}}
</consulta>
"""

prompt_v1 = mlflow.genai.register_prompt(
    name=PROMPT_NAME,
    template=template_v1,
    commit_message="v1: plantilla baseline con delimitadores XML",
    tags={"dominio": "soporte", "idioma": "es", "autor": "equipo-llmops"},
)
print(f"Registrado: {prompt_v1.uri} (version {prompt_v1.version})")

# 2. Iteracion a v2: politica de abstencion y salida estructurada JSON
template_v2 = """Eres un ingeniero de soporte senior. Responde UNICAMENTE con hechos
verificables en <contexto>. Si la evidencia es insuficiente, establece
"evidencia_suficiente": false y solicita los datos faltantes sin inventar nada.

Devuelve exclusivamente un JSON valido con este esquema:
{"respuesta": "...", "evidencia_suficiente": true, "citas": ["DOC-1"]}

<contexto>
{{context}}
</contexto>

<consulta>
{{question}}
</consulta>
"""

prompt_v2 = mlflow.genai.register_prompt(
    name=PROMPT_NAME,
    template=template_v2,
    commit_message="v2: anade abstencion explicita, citas y contrato JSON",
    tags={"dominio": "soporte", "idioma": "es", "formato": "json"},
)

# 3. Gestion de aliases mutables (promocion gobernada)
mlflow.genai.set_prompt_alias(
    name=PROMPT_NAME,
    alias="staging",
    version=prompt_v2.version,
)
# Tras pasar las puertas de calidad, se promueve a @production
mlflow.genai.set_prompt_alias(
    name=PROMPT_NAME,
    alias="production",
    version=prompt_v2.version,
)

# Inspeccion de metadatos y versiones a traves de MlflowClient
client = MlflowClient()
prod_meta = client.load_prompt(f"prompts:/{PROMPT_NAME}@production")
print(
    f"Alias @production apunta a v{prod_meta.version} | "
    f"Commit: {prod_meta.commit_message} | Variables: {prod_meta.variables}"
)

# 4. Carga dinamica en el microservicio de produccion con cache TTL
# Evita realizar una peticion HTTP al servidor MLflow en cada consulta de usuario
def responder_consulta_produccion(question: str, context: str) -> str:
    prompt = mlflow.genai.load_prompt(
        f"prompts:/{PROMPT_NAME}@production",
        cache_ttl_seconds=300,  # Refresco automatico cada 5 minutos
    )
    # Formateo seguro: las llaves simples del JSON no interfieren con {{var}}
    rendered_prompt = prompt.format(question=question, context=context)

    openai_client = OpenAI()
    resp = openai_client.chat.completions.create(
        model="gpt-4o-mini-2024-07-18",
        temperature=0.0,
        messages=[{"role": "user", "content": rendered_prompt}],
    )
    return resp.choices[0].message.content

# 5. Vinculacion de linaje entre LoggedModel, version del prompt y evaluacion
with mlflow.start_run(run_name="registro_linaje_v2"):
    logged_model = mlflow.set_active_model(name="agente_soporte_rag")
    # Asociar explicitamente la URI inmutable del prompt al modelo registrado
    mlflow.log_params({
        "prompt_uri_inmutable": prompt_v2.uri,  # prompts:/soporte_cliente/2
        "prompt_alias_origen": f"prompts:/{PROMPT_NAME}@production",
        "logged_model_id": logged_model.model_id,
    })
    mlflow.set_tag("mlflow.linkedPrompts", prompt_v2.uri)
\end{lstlisting}

Dos detalles de arquitectura en el Listado~\ref{lst:prompt-registry-workflow} merecen especial atención en producción:
\begin{enumerate}
  \item \textbf{Resiliencia y caché de cliente:} habilitar \code{cache\_ttl\_seconds} (por ejemplo 300 segundos) evita introducir un punto único de fallo y elimina la latencia de red hacia el servidor de MLflow en la ruta crítica de cada petición de usuario. Como buena práctica defensiva adicional, conviene mantener una plantilla de respaldo empaquetada o precargada durante el arranque del contenedor por si el registro se encontrase temporalmente inaccesible en el primer inicio.
  \item \textbf{Alias para resolución, versión numérica para linaje:} la aplicación carga \code{@production} para beneficiarse del desacoplamiento, pero las trazas y los \code{LoggedModel} registran siempre el número de versión inmutable resuelto (por ejemplo \code{prompts:/soporte\_cliente/2}). De este modo, si una semana después \code{@production} apunta a la versión \code{5}, el historial de auditoría conserva con exactitud qué versión concreta generó cada respuesta histórica.
\end{enumerate}

% ============================================================================
\section[Optimización Automática: DSPy y GEPA]{Optimización Automática de Prompts: DSPy y GEPA con MLflow}\label{sec:prompt-optimization}
% ============================================================================

El versionado en un registro resuelve el gobierno de los \textit{prompts}, pero no elimina el cuello de botella creativo y estadístico de redactarlos a mano. Conforme las aplicaciones GenAI crecen en complejidad ---integrando recuperación híbrida, múltiples herramientas y decenas de reglas de negocio--- la afinación manual alcanza rápidamente un techo asintótico.

En primer lugar, el espacio de búsqueda de instrucciones en lenguaje natural es prácticamente infinito y no admite intuición fiable: reformulaciones que parecen más claras a ojos humanos a menudo empeoran el rendimiento en un modelo concreto, mientras que pequeñas restricciones estructurales pueden elevar sustancialmente la precisión. En segundo lugar, los \textit{prompts} escritos a mano sufren de \textbf{sobreajuste al modelo base}: cuando una organización migra de \code{gpt-4o} a \code{claude-3-7-sonnet}, a \code{gemini-2.5-pro} o a un modelo abierto servido con \pkg{vLLM} para reducir costes, semanas de ajustes manuales pierden vigencia de la noche a la mañana.

Entre 2024 y 2026 ha madurado una respuesta ingenieril rigurosa a este problema: \textbf{tratar los \textit{prompts} no como texto escrito manualmente, sino como parámetros discretos compilados automáticamente frente a una métrica de evaluación sobre datos de entrenamiento}. En el ecosistema moderno integrado con \pkg{MLflow 3} destacan dos enfoques complementarios: la programación declarativa de \textbf{DSPy} y la optimización reflexiva evolutiva \textbf{GEPA} integrada nativamente en \code{mlflow.genai.optimize\_prompts()}.

\subsection{Programación declarativa y compilación de prompts con DSPy}\label{subsec:dspy-mlflow}

\textbf{DSPy} (\textit{Declarative Self-improving Python}, desarrollado en la Universidad de Stanford) propone separar radicalmente \emph{qué} debe transformar un flujo LLM de \emph{cómo} se redacta el \textit{prompt} concreto para lograrlo. Su arquitectura descansa sobre tres abstracciones fundamentales:

\begin{itemize}
  \item \textbf{Signatures (Firmas):} especifican de forma declarativa y tipada los campos de entrada y salida de un paso LLM (por ejemplo, \code{contexto, consulta -> diagnostico, cita\_documento}), actuando como contrato semántico independiente del modelo subyacente.
  \item \textbf{Modules (Módulos):} bloques parametrizados que implementan estrategias de inferencia sobre una firma, como \code{dspy.Predict}, \code{dspy.ChainOfThought} o \code{dspy.ReAct}. Pueden componerse en grafos complejos de manera análoga a las capas de \pkg{PyTorch}.
  \item \textbf{Optimizers (antes \textit{Teleprompters}):} compiladores automáticos que, dada una métrica de evaluación escalar $\mu(y, \hat{y})$ y un conjunto de ejemplos de entrenamiento, optimizan las instrucciones textuales y las demostraciones \textit{few-shot} de cada módulo. Entre los más utilizados destacan \code{BootstrapFewShot} (genera trazas válidas con un modelo profesor y selecciona las demostraciones que maximizan $\mu$) y \code{MIPROv2} (\textit{Multiprompt Instruction Proposal Optimizer}, que propone instrucciones candidatas informadas por los datos y utiliza optimización bayesiana sobre pares instrucción--demostraciones).
\end{itemize}

La integración nativa \code{mlflow.dspy.autolog()} convierte el proceso de compilación de DSPy ---a menudo opaco si se ejecuta aislado--- en una experiencia completamente observable: registra cada ensayo intermedio del optimizador como un \textit{run} anidado, captura las trazas de razonamiento durante la compilación, almacena el programa DSPy compilado como artefacto versionado y grafica la curva de convergencia de la métrica. El Listado~\ref{lst:dspy-mlflow} ilustra este patrón completo.

\begin{lstlisting}[style=python, caption={Compilación automática de un módulo RAG/triaje con DSPy (\code{MIPROv2}) y seguimiento completo mediante \code{mlflow.dspy.autolog()}.}, label={lst:dspy-mlflow}]
import dspy
import mlflow

# 1. Activar autologging completo de DSPy en MLflow 3
mlflow.set_experiment("optimizacion-dspy-soporte")
mlflow.dspy.autolog(
    log_compiles=True,              # Registra cada iteracion del optimizador
    log_evals=True,                 # Registra metricas sobre el conjunto de validacion
    log_traces_from_compile=True,   # Guarda trazas generadas durante la compilacion
)

# 2. Configurar el modelo objetivo (rapido y economico para produccion)
lm = dspy.LM(model="openai/gpt-4o-mini-2024-07-18", temperature=0.0)
dspy.configure(lm=lm)

# 3. Declarar Signature y Module de forma independiente del prompt textual
class ClasificarSoporte(dspy.Signature):
    """Clasifica un ticket de soporte y justifica si requiere escalado humano."""
    politica_empresa: str = dspy.InputField(desc="Extracto vigente de politicas SLA")
    ticket_usuario: str = dspy.InputField(desc="Mensaje original del cliente")
    categoria: str = dspy.OutputField(desc="facturacion | tecnico | envio | cancelacion")
    requiere_escalado: bool = dspy.OutputField(desc="true si supera umbral SLA o riesgo")

class TriajeModule(dspy.Module):
    def __init__(self):
        super().__init__()
        self.clasificador = dspy.ChainOfThought(ClasificarSoporte)

    def forward(self, politica_empresa: str, ticket_usuario: str):
        return self.clasificador(
            politica_empresa=politica_empresa,
            ticket_usuario=ticket_usuario,
        )

# 4. Metrica de evaluacion para guiar al compilador
def metrica_exactitud(example, pred, trace=None) -> float:
    cat_ok = pred.categoria.strip().lower() == example.categoria.strip().lower()
    esc_ok = bool(pred.requiere_escalado) == bool(example.requiere_escalado)
    return 1.0 if (cat_ok and esc_ok) else 0.0

# Ejemplos ilustrativos (en produccion usar 50-300 ejemplos curados)
trainset = [
    dspy.Example(
        politica_empresa="Cobros duplicados y bajas corporativas escalan a Nivel 2.",
        ticket_usuario="Me han cobrado dos veces la suscripcion este mes.",
        categoria="facturacion",
        requiere_escalado=True,
    ).with_inputs("politica_empresa", "ticket_usuario"),
    dspy.Example(
        politica_empresa="Errores de interfaz sin perdida de datos se resuelven en Nivel 1.",
        ticket_usuario="No veo el boton de descarga de facturas en modo oscuro.",
        categoria="tecnico",
        requiere_escalado=False,
    ).with_inputs("politica_empresa", "ticket_usuario"),
]

# 5. Compilacion con MIPROv2 bajo un run de MLflow
with mlflow.start_run(run_name="compilacion_miprov2_gpt4o_mini"):
    optimizer = dspy.MIPROv2(
        metric=metrica_exactitud,
        auto="medium",              # Presupuesto equilibrado de ensayos bayesianos
        num_threads=4,
    )
    compiled_module = optimizer.compile(
        TriajeModule(),
        trainset=trainset,
        max_bootstrapped_demos=3,
        max_labeled_demos=3,
    )

    # Registrar el programa DSPy compilado como modelo reproducible
    mlflow.dspy.log_model(
        dspy_model=compiled_module,
        name="triaje_dspy_compilado",
        input_example={
            "politica_empresa": "Cobros duplicados escalan a Nivel 2.",
            "ticket_usuario": "Cargo duplicado en tarjeta corporativa.",
        },
    )
\end{lstlisting}

\subsection{Optimización Reflexiva y Evolutiva (GEPA) nativa en MLflow 3}\label{subsec:gepa-mlflow}

Mientras que DSPy exige adoptar su modelo de programación declarativo, muchas organizaciones disponen ya de aplicaciones en producción escritas en Python nativo, LangGraph, LlamaIndex o el SDK de OpenAI cuyos \textit{prompts} residen en el MLflow Prompt Registry. Para ese escenario agnóstico de \textit{framework}, \pkg{MLflow 3} incorpora nativamente la función \code{mlflow.genai.optimize\_prompts()} respaldada por el algoritmo \textbf{GEPA} (\textit{Genetic-Pareto / Reflective Prompt Adaptation}).

A diferencia de los optimizadores puramente escalares ---que solo observan si una respuesta obtuvo $0$ o $1$ y buscan casi a ciegas---, \textbf{GEPA explota la riqueza diagnóstica del lenguaje natural}:

\begin{enumerate}
  \item \textbf{Evaluación con trazas y justificaciones:} el optimizador ejecuta la aplicación (\code{predict\_fn}) sobre un minilote del conjunto de entrenamiento usando el \textit{prompt} candidato actual. Los \textit{scorers} (tanto jueces LLM como validadores deterministas en Python) no solo devuelven una puntuación, sino una explicación textual (\textit{rationale}) del motivo exacto del fallo y acceso a la traza completa de ejecución.
  \item \textbf{Reflexión guiada por un modelo profesor (\textit{reflection\_model}):} GEPA entrega las trazas fallidas, el \textit{prompt} actual y los diagnósticos de los jueces a un LLM de alta capacidad (por ejemplo \code{gpt-4o} o \code{claude-sonnet}) que actúa como ingeniero revisor: identifica patrones sistemáticos de error (instrucciones ambiguas, ausencia de reglas de prioridad, formatos mal especificados) y sintetiza una mutación dirigida de la plantilla.
  \item \textbf{Archivo evolutivo de Pareto:} en lugar de mantener una única rama lineal que podría sobreajustarse al último minilote corregido, GEPA conserva un conjunto de candidatos no dominados a lo largo de los distintos casos de evaluación (frente de Pareto por instancia) y combina rasgos ganadores entre candidatos complementarios.
  \item \textbf{Registro automático con linaje completo:} al converger o agotarse el presupuesto de llamadas, \code{mlflow.genai.optimize\_prompts()} registra la nueva versión optimizada directamente en el \textbf{MLflow Prompt Registry}, preservando el linaje exacto entre la versión semilla, el historial de iteraciones y la ganancia medida en cada \textit{scorer}.
\end{enumerate}

El Listado~\ref{lst:gepa-mlflow} demuestra cómo aplicar \code{GepaPromptOptimizer} directamente sobre un \textit{prompt} previamente registrado en el Prompt Registry, utilizando un modelo frontera como motor de reflexión para optimizar un \textit{prompt} que será ejecutado en producción por un modelo ligero y diez veces más barato.

\begin{lstlisting}[style=python, caption={Optimización automática de un prompt registrado mediante GEPA en MLflow 3 (\code{mlflow.genai.optimize\_prompts}).}, label={lst:gepa-mlflow}]
import json
import mlflow
from mlflow.entities import Feedback
from mlflow.genai.optimize.optimizers import GepaPromptOptimizer
from mlflow.genai.scorers import Correctness, scorer
from openai import OpenAI

mlflow.set_experiment("optimizacion-gepa-soporte")
openai_client = OpenAI()

PROMPT_NAME = "soporte_cliente"

# 1. Registrar plantilla semilla sencilla (v1) en Prompt Registry
seed_prompt = mlflow.genai.register_prompt(
    name=PROMPT_NAME,
    template=(
        "Responde a la consulta usando el contexto y devuelve un JSON con "
        "las claves respuesta y escalar.\n"
        "Contexto: {{context}}\nConsulta: {{question}}"
    ),
    commit_message="v1: semilla inicial antes de optimizacion GEPA",
)

# 2. Scorer determinista en Python cuyo rationale alimenta la reflexion de GEPA
@scorer
def json_y_politica_escalado(outputs: str) -> Feedback:
    try:
        data = json.loads(outputs)
    except json.JSONDecodeError as exc:
        return Feedback(
            value=False,
            rationale=f"La salida no es JSON valido (error: {exc}). Debe omitir markdown.",
        )
    required = {"respuesta", "escalar"}
    missing = required - set(data.keys())
    if missing:
        return Feedback(
            value=False,
            rationale=f"Faltan claves obligatorias en el JSON: {sorted(missing)}.",
        )
    if not isinstance(data["escalar"], bool):
        return Feedback(
            value=False,
            rationale="El campo 'escalar' debe ser booleano (true/false), no cadena.",
        )
    return Feedback(value=True, rationale="Contrato JSON estrictamente cumplido.")

# 3. Funcion de prediccion que carga dinamicamente la plantilla evaluada por GEPA
def predict_fn(context: str, question: str) -> str:
    # Durante optimize_prompts, MLflow inyecta la variante candidata sobre esta URI
    prompt = mlflow.genai.load_prompt(seed_prompt.uri)
    rendered = prompt.format(context=context, question=question)
    response = openai_client.chat.completions.create(
        model="gpt-4o-mini-2024-07-18",
        temperature=0.0,
        messages=[{"role": "user", "content": rendered}],
    )
    return response.choices[0].message.content

# 4. Conjunto de entrenamiento con expectativas de referencia
train_data = [
    {
        "inputs": {
            "context": "Politica: los reembolsos superiores a 500 EUR requieren aprobacion humana.",
            "question": "Pedi un reembolso de 620 EUR por un equipo defectuoso.",
        },
        "expectations": {
            "expected_response": (
                '{"respuesta": "Su reembolso de 620 EUR supera el umbral de 500 EUR y '
                'ha sido derivado a un agente para aprobacion.", "escalar": true}'
            )
        },
    },
    {
        "inputs": {
            "context": "Politica: los cambios de direccion antes del envio son automaticos.",
            "question": "Todavia no ha salido mi pedido, puedo cambiar la calle de entrega?",
        },
        "expectations": {
            "expected_response": (
                '{"respuesta": "Si, dado que el pedido aun no se ha enviado, puede '
                'actualizar la direccion automaticamente desde su panel.", "escalar": false}'
            )
        },
    },
    {
        "inputs": {
            "context": "Politica: el soporte cubre exclusivamente productos adquiridos en tienda oficial.",
            "question": "Compre el dispositivo de segunda mano y falla la bateria, me lo cambian?",
        },
        "expectations": {
            "expected_response": (
                '{"respuesta": "La garantia de soporte cubre unicamente compras realizadas '
                'en la tienda oficial, por lo que no aplica sustitucion.", "escalar": false}'
            )
        },
    },
]

# 5. Ejecutar optimizacion reflexiva GEPA (profesor: gpt-4o -> alumno: gpt-4o-mini)
opt_result = mlflow.genai.optimize_prompts(
    predict_fn=predict_fn,
    train_data=train_data,
    prompt_uris=[seed_prompt.uri],
    optimizer=GepaPromptOptimizer(
        reflection_model="openai:/gpt-4o",
        max_metric_calls=60,
    ),
    scorers=[
        Correctness(model="openai:/gpt-4o"),
        json_y_politica_escalado,
    ],
    enable_tracking=True,
)

# 6. Inspeccionar la version ganadora registrada automaticamente y asignar @staging
best_prompt = opt_result.optimized_prompts[0]
print(f"Prompt optimizado registrado en: {best_prompt.uri}")
print(f"Mejora obtenida: {opt_result.initial_eval_score:.3f} -> {opt_result.final_eval_score:.3f}")

mlflow.genai.set_prompt_alias(
    name=PROMPT_NAME,
    alias="staging",
    version=best_prompt.version,
)
\end{lstlisting}

\subsection{Análisis de la Frontera de Pareto: Calidad frente a Coste y Latencia}\label{subsec:frontera-pareto}

Una de las conclusiones más valiosas ---y financieramente rentables--- de aplicar optimización sistemática de \textit{prompts} en producción aparece al analizar el espacio multiobjetivo de calidad y economía. Sea $\mathcal{S}$ el conjunto de configuraciones candidatas $s = (m, p)$, donde $m$ denota el modelo fundacional y $p$ la plantilla de \textit{prompt}. Cada configuración induce un vector de desempeño $(Q(s), C(s), L_{95}(s))$, representando respectivamente calidad compuesta en el conjunto de evaluación (a maximizar), coste por cada 1\,000 peticiones (a minimizar) y latencia p95 (a minimizar).

Decimos que una configuración $s_1$ \textbf{domina en sentido de Pareto} a $s_2$ ($s_1 \succ s_2$) si y solo si no es peor en ninguna dimensión y es estrictamente mejor al menos en una:
\begin{equation}\label{eq:pareto-dominancia}
\begin{aligned}
  s_1 \succ s_2 \iff \;& \Big(Q(s_1) \ge Q(s_2) \;\wedge\; C(s_1) \le C(s_2) \;\wedge\; L_{95}(s_1) \le L_{95}(s_2)\Big) \\
  &\wedge\; \exists\, d \in \{Q, C, L_{95}\} \text{ con desigualdad estricta.}
\end{aligned}
\end{equation}
El subconjunto de configuraciones no dominadas conforma la \textbf{Frontera de Pareto}. Ninguna organización racional debería desplegar en producción una solución estrictamente dominada situada por debajo y a la derecha de dicha frontera.

La Figura~\ref{fig:pareto-prompt-opt} ilustra un fenómeno observado reiteradamente en cargas empresariales reales (clasificación, extracción estructurada, enrutamiento y RAG focalizado): un modelo ligero como \code{gpt-4o-mini} con un \textit{prompt} \textit{zero-shot} artesanal obtiene una exactitud modesta (74\,\%), lo que frecuentemente empuja a los equipos a saltar directamente a un modelo frontera caro como \code{gpt-4o} (88\,\%--93\,\% a un coste más de $12\times$ superior). Sin embargo, al aplicar compilación automática con \textbf{GEPA} o \textbf{DSPy} utilizando el modelo frontera únicamente como profesor fuera de línea (\textit{offline}), el modelo ligero asciende hasta el 91.2\,\% de exactitud con un coste inferior a \$0.35 por cada 1\,000 peticiones ---desplazando radicalmente la frontera de Pareto y dominando incluso a la configuración \textit{zero-shot} del modelo frontera.

\begin{figure}[htbp]
\centering
\begin{tikzpicture}
\begin{axis}[
    width=0.94\textwidth,
    height=8.4cm,
    xlabel={\sffamily\small Coste estimado por 1\,000 peticiones (USD, escala logarítmica)},
    ylabel={\sffamily\small Exactitud compuesta en Golden Dataset (\%)},
    xmode=log,
    xmin=0.12, xmax=18,
    ymin=68, ymax=102,
    xtick={0.2, 0.5, 1, 2, 5, 10},
    xticklabels={\$0.20, \$0.50, \$1.00, \$2.00, \$5.00, \$10.00},
    ytick={70, 75, 80, 85, 90, 95, 100},
    grid=both,
    major grid style={LightBorder, line width=0.5pt},
    minor grid style={LightBorder!45, dashed, line width=0.3pt},
    tick label style={font=\sffamily\footnotesize, color=SlateText},
    axis line style={PrimaryDark, line width=0.8pt},
    legend pos=south east,
    legend style={
      font=\sffamily\scriptsize,
      fill=white,
      draw=LightBorder,
      rounded corners=2pt,
      cells={anchor=west},
      inner sep=4pt
    }
]

% Curva de la Frontera de Pareto
\addplot[
    color=TealAccent,
    line width=1.6pt,
    dashed,
    mark=none
] coordinates {
    (0.18, 74.0)
    (0.32, 91.2)
    (1.45, 93.4)
    (8.60, 95.1)
};
\addlegendentry{Frontera de Pareto (óptimo calidad--coste)}

% Puntos en la frontera (no dominados)
\addplot[
    only marks,
    mark=*,
    mark size=3.3pt,
    color=TealAccent
] coordinates {
    (0.18, 74.0)
    (0.32, 91.2)
    (1.45, 93.4)
    (8.60, 95.1)
};
\addlegendentry{Configuraciones Pareto-óptimas}

% Puntos dominados (suboptimos)
\addplot[
    only marks,
    mark=square*,
    mark size=3.0pt,
    color=AmberAccent
] coordinates {
    (0.48, 82.5)
    (3.90, 88.2)
    (5.80, 90.4)
};
\addlegendentry{Configuraciones dominadas (subóptimas)}

% Anotaciones limpias con punteros donde aporta claridad (cero solapamientos)
\node[font=\sffamily\scriptsize, color=PrimaryDark, anchor=west, align=left]
  at (axis cs:0.21, 73.6) {Zero-Shot\\[-1pt]modelo ligero};

\node[font=\sffamily\scriptsize, color=AmberAccent, anchor=west, align=left]
  at (axis cs:0.55, 82.0) {Few-Shot manual\\[-1pt]modelo ligero};

\node[font=\sffamily\scriptsize\bfseries, color=TealAccent, anchor=south west, align=left,
      fill=ModernBg, draw=TealAccent, rounded corners=2pt, inner sep=3pt] (gepa_mini)
  at (axis cs:0.15, 95.8) {GEPA / DSPy en modelo ligero\\[-1pt](91.2\% por \$0.32/1k --- ahorro $>12\times$)};
\draw[->, >=Stealth, line width=0.8pt, color=TealAccent]
  (gepa_mini.south) -- (axis cs:0.32, 91.9);

\node[font=\sffamily\scriptsize, color=PrimaryDark, anchor=south, align=center] (gepa_mid)
  at (axis cs:1.85, 96.4) {GEPA modelo\\[-1pt]intermedio (93.4\%)};
\draw[->, >=Stealth, line width=0.7pt, color=MutedGray]
  (gepa_mid.south) -- (axis cs:1.48, 94.0);

\node[font=\sffamily\scriptsize, color=PrimaryDark, anchor=south, align=center]
  at (axis cs:8.60, 96.3) {CoT + GEPA\\[-1pt]modelo frontera (95.1\%)};

\node[font=\sffamily\scriptsize, color=AmberAccent, anchor=east, align=right]
  at (axis cs:3.45, 87.7) {Zero-Shot modelo\\[-1pt]frontera (dominado)};

\node[font=\sffamily\scriptsize, color=AmberAccent, anchor=west, align=left]
  at (axis cs:6.50, 89.8) {Few-Shot manual\\[-1pt]modelo frontera};

\end{axis}
\end{tikzpicture}
\caption{Frontera de Pareto (Calidad frente a Coste por 1\,000 peticiones en escala logarítmica) para distintas combinaciones de modelo fundacional y técnica de \textit{prompting}. La optimización automática con GEPA/DSPy sobre un modelo ligero domina estrictamente al enfoque \textit{zero-shot} e incluso al \textit{few-shot} manual sobre un modelo frontera mucho más costoso.}
\label{fig:pareto-prompt-opt}
\end{figure}

Este patrón constituye una forma práctica de \textbf{destilación de conocimiento sin reentrenamiento de pesos}: se invierte presupuesto de cómputo una sola vez en un modelo profesor potente durante la fase \textit{offline} de optimización, y el beneficio se amortiza millones de veces durante la inferencia productiva en un modelo rápido y económico.

% ============================================================================
\section[Golden Datasets y CI/CD para Prompts]{Golden Datasets, Pruebas de Regresión y CI/CD para Prompts}\label{sec:golden-datasets-cicd}
% ============================================================================

Tanto un barrido manual como un optimizador avanzado como GEPA son tan fiables como el conjunto de evaluación contra el que miden el progreso. Si las pruebas se limitan a diez ejemplos limpios redactados por el propio desarrollador, cualquier optimizador tenderá a sobreajustarse y fallará ante el ruido del tráfico real. Por ello, el último pilar del \textit{prompt engineering} profesional consiste en construir un \textbf{Golden Evaluation Dataset} versionado e integrarlo en una puerta de regresión automatizada dentro del flujo CI/CD.

\subsection{Diseño y estratificación de un Golden Evaluation Dataset}\label{subsec:golden-dataset}

Un \textit{Golden Dataset} maduro (típicamente entre 100 y 500 casos cuidadosamente revisados por expertos de dominio, crecido de forma orgánica desde las trazas reales de producción como veremos en la Sección~\ref{sec:llmops-flywheel}) no debe ser una muestra aleatoria ingenua, sino una muestra \textbf{estratificada por perfiles de riesgo}:

\begin{enumerate}
  \item \textbf{Casos canónicos de distribución principal (50--60\,\%):} reflejan las intenciones, estilos de redacción, faltas ortográficas y longitudes más frecuentes en producción. Garantizan que ninguna optimización degrada la experiencia del usuario mayoritario.
  \item \textbf{Casos borde de dominio y formato (15--20\,\%):} consultas con múltiples intenciones simultáneas, cifras ambiguas, tablas incrustadas, cambios de idioma o reglas de negocio que se solapan en la frontera de decisión.
  \item \textbf{Consultas subespecificadas o ambiguas (10\,\%):} peticiones en las que falta un parámetro esencial (por ejemplo, número de pedido o región fiscal) y el comportamiento correcto consiste en formular una pregunta aclaratoria en lugar de adivinar.
  \item \textbf{Casos negativos o \textit{out-of-scope} (10\,\%):} preguntas ajenas al dominio corporativo o carentes de soporte documental, diseñadas para medir la tasa de abstención honesta y penalizar alucinaciones.
  \item \textbf{Casos adversariales y de seguridad (5--10\,\%):} intentos explícitos de \textit{prompt injection} (directo e indirecto dentro de documentos recuperados), \textit{jailbreak}, extracción del \textit{system prompt} y solicitudes que inducen la exposición de datos personales (PII). En este estrato la tolerancia a regresión suele ser estrictamente cero.
\end{enumerate}

En \pkg{MLflow 3}, estos conjuntos se versionan como artefactos evaluables mediante \code{mlflow.genai.datasets} (o \code{mlflow.data.from\_pandas} con digest criptográfico automático), garantizando que el historial de pruebas pueda responder en todo momento qué versión exacta del \textit{Golden Dataset} certificó cada promoción de \textit{prompt}. Adicionalmente, cuando se emplean optimizadores como DSPy o GEPA, resulta imprescindible mantener una \textbf{partición estricta entre \code{train\_set} de optimización y \code{holdout\_golden\_set} de certificación}: evaluar la puerta de promoción sobre los mismos ejemplos vistos por el optimizador produce estimaciones artificialmente optimistas análogas a la fuga de datos (\textit{data leakage}) en ML clásico.

\subsection{Puertas de regresión (\textit{Regression Gating}) y promoción CI/CD}\label{subsec:regression-gating}

En un flujo CI/CD moderno para \textit{prompts} ---ejecutado desde GitHub Actions, GitLab CI o un DAG de Apache Airflow--- ninguna modificación de plantilla promueve directamente al alias \code{@production}. En su lugar, la nueva versión recibe el alias temporal \code{@candidate} (o \code{@staging}) y dispara un trabajo automatizado que ejecuta una comparación pareada (\textit{paired evaluation}) frente al campeón actual \code{@production} sobre el mismo \textit{Golden Dataset} congelado.

Una política robusta de \textit{regression gating} combina cuatro tipos de restricciones simultáneas:
\begin{itemize}
  \item \textbf{Restricción de calidad mínima y no regresión:} la exactitud o puntuación compuesta del candidato debe superar un umbral absoluto mínimo ($Q_{\text{cand}} \ge Q_{\min}$) y no degradar respecto al campeón más allá de un margen de tolerancia estadística ($Q_{\text{cand}} - Q_{\text{prod}} \ge -\delta$, o exigir mejora estricta $\Delta Q > 0$ mediante una prueba pareada de McNemar o \textit{bootstrap}).
  \item \textbf{Guardrail de seguridad de tolerancia cero:} la tasa de aprobación en el estrato adversarial y de seguridad (\code{safety\_pass\_rate}) debe ser del 100\,\% (o estrictamente no inferior al 99.5\,\% en baterías extensas), bloqueando de inmediato cualquier fuga de instrucciones o respuesta tóxica.
  \item \textbf{Techo de alucinación:} la tasa de fidelidad al contexto (\code{groundedness}) no puede caer por debajo del SLA acordado (por ejemplo $\ge 0.95$).
  \item \textbf{Presupuesto operacional y económico:} la latencia p95 y el coste estimado por 1\,000 peticiones no deben superar los techos definidos por FinOps (por ejemplo $L_{95} \le 2\,500\text{ ms}$ y $C_{1\text{k}} \le \$0.60$).
\end{itemize}

El Listado~\ref{lst:prompt-regression-gate} implementa un script de puerta de regresión invocable directamente desde un \textit{pipeline} CI/CD. Si todas las restricciones se satisfacen, el script promueve atómicamente el alias \code{@production} en el MLflow Prompt Registry; en caso contrario, finaliza con código de salida distinto de cero para bloquear la fusión o el despliegue, dejando el informe detallado en MLflow Tracking.

\begin{lstlisting}[style=python, caption={Script de puerta de regresión CI/CD comparando \code{@candidate} frente a \code{@production} antes de promover el alias en MLflow Prompt Registry.}, label={lst:prompt-regression-gate}]
import sys
import mlflow
import pandas as pd
from mlflow.genai.scorers import Correctness, Safety
from openai import OpenAI

PROMPT_NAME = "soporte_cliente"
GOLDEN_DATASET_PATH = "datasets/golden_soporte_v3.jsonl"
MODEL_SNAPSHOT = "gpt-4o-mini-2024-07-18"

# Politica de Regression Gating (thresholds de negocio y FinOps)
GATE_POLICY = {
    "min_correctness": 0.90,
    "max_correctness_drop": 0.005,   # Tolerancia maxima frente a @production
    "min_safety_rate": 1.00,         # Tolerancia cero en casos adversariales
    "max_cost_per_1k_usd": 0.55,
    "max_latency_p95_ms": 2200.0,
}

client = OpenAI()

def build_predict_fn(prompt_uri: str):
    prompt = mlflow.genai.load_prompt(prompt_uri)

    def _predict(context: str, question: str) -> str:
        rendered = prompt.format(context=context, question=question)
        resp = client.chat.completions.create(
            model=MODEL_SNAPSHOT,
            temperature=0.0,
            messages=[{"role": "user", "content": rendered}],
        )
        return resp.choices[0].message.content

    return _predict

def evaluate_alias(alias: str, eval_df: pd.DataFrame) -> dict:
    uri = f"prompts:/{PROMPT_NAME}@{alias}"
    prompt_obj = mlflow.genai.load_prompt(uri)
    with mlflow.start_run(run_name=f"gate_{alias}_v{prompt_obj.version}"):
        mlflow.log_params({
            "prompt_alias": alias,
            "prompt_version": prompt_obj.version,
            "golden_dataset": GOLDEN_DATASET_PATH,
            "model_snapshot": MODEL_SNAPSHOT,
        })
        # Registrar linaje formal del Golden Dataset en MLflow
        dataset = mlflow.data.from_pandas(
            eval_df, source=GOLDEN_DATASET_PATH, name="golden_soporte_v3"
        )
        mlflow.log_input(dataset, context="regression_gate")

        res = mlflow.genai.evaluate(
            data=eval_df,
            predict_fn=build_predict_fn(uri),
            scorers=[Correctness(model="openai:/gpt-4o"), Safety()],
        )
        return {
            "version": prompt_obj.version,
            "correctness": res.metrics.get("correctness/mean", 0.0),
            "safety": res.metrics.get("safety/mean", 0.0),
            "latency_p95": res.metrics.get("latency_ms/p95", 0.0),
            "cost_1k": res.metrics.get("cost_per_1k_usd", 0.0),
        }

def main() -> None:
    mlflow.set_experiment("ci-cd-regression-gate-prompts")
    eval_df = pd.read_json(GOLDEN_DATASET_PATH, lines=True)

    prod_stats = evaluate_alias("production", eval_df)
    cand_stats = evaluate_alias("candidate", eval_df)

    failures = []
    if cand_stats["correctness"] < GATE_POLICY["min_correctness"]:
        failures.append(
            f"Exactitud insuficiente: {cand_stats['correctness']:.3f} "
            f"< {GATE_POLICY['min_correctness']:.3f}"
        )
    if (prod_stats["correctness"] - cand_stats["correctness"]) > GATE_POLICY["max_correctness_drop"]:
        failures.append(
            f"Regresion frente a @production: {cand_stats['correctness']:.3f} "
            f"vs {prod_stats['correctness']:.3f}"
        )
    if cand_stats["safety"] < GATE_POLICY["min_safety_rate"]:
        failures.append(f"Fallo de seguridad adversarial: {cand_stats['safety']:.3f} < 1.000")
    if cand_stats["cost_1k"] > GATE_POLICY["max_cost_per_1k_usd"]:
        failures.append(f"Presupuesto FinOps excedido: ${cand_stats['cost_1k']:.3f}/1k")
    if cand_stats["latency_p95"] > GATE_POLICY["max_latency_p95_ms"]:
        failures.append(f"Latencia p95 excedida: {cand_stats['latency_p95']:.1f} ms")

    if failures:
        print("BLOQUEO DE PROMOCION (@candidate rechazado):")
        for msg in failures:
            print(f"  - {msg}")
        sys.exit(1)

    # Promocion atomica del alias @production a la version verificada
    mlflow.genai.set_prompt_alias(
        name=PROMPT_NAME,
        alias="production",
        version=cand_stats["version"],
    )
    print(
        f"PROMOCION APROBADA: {PROMPT_NAME}@production -> v{cand_stats['version']} "
        f"(correctness={cand_stats['correctness']:.3f}, safety={cand_stats['safety']:.3f})"
    )

if __name__ == "__main__":
    main()
\end{lstlisting}

\begin{infobox}[De la experimentación de prompts a sistemas GenAI completos]
En este capítulo hemos establecido los cimientos del LLMOps moderno: cómo experimentar con \textit{prompts} de forma reproducible en MLflow Tracking, cómo versionarlos y promocionarlos sin acoplamiento mediante el Prompt Registry, cómo compilarlos automáticamente con DSPy y GEPA para desplazar la frontera de Pareto y cómo blindar cada cambio con \textit{Golden Datasets} y puertas de regresión. Sin embargo, las aplicaciones generativas de producción rara vez consisten en una única llamada aislada a un LLM: combinan recuperadores vectoriales, \textit{rerankers}, invocación iterativa de herramientas, múltiples proveedores y millones de interacciones vivas. En el Capítulo~\ref{chap:mlflow-llms} abordamos precisamente esa capa completa: observabilidad distribuida con \textbf{MLflow Tracing}, evaluación avanzada de sistemas RAG y agentes con \code{mlflow.genai.evaluate}, empaquetado sin \code{pickle} mediante \code{ResponsesAgent}, gobernanza centralizada con el \textbf{MLflow AI Gateway} y monitorización continua en producción.
\end{infobox}
